% >>> j2l compat: xcolor / named-colour compatibility
\PassOptionsToPackage{dvipsnames}{color}
\PassOptionsToPackage{dvipsnames,svgnames,x11names,table}{xcolor}
\RequirePackage{xcolor}
\makeatletter\@namedef{ver@color.sty}{2021/01/01}\makeatother
% <<< j2l compat
\documentclass{default}
\IfFileExists{default.sty}{\IfFileExists{default.sty}{\usepackage{default}}{}}{}

\IfFileExists{float.sty}{\IfFileExists{float.sty}{\usepackage{float}}{}}{}
% J2L-OPT-BEGIN (auto-generated; regenerated every iteration)
\makeatletter
\@ifpackageloaded{geometry}{}{\usepackage{geometry}}
\geometry{margin=1.00in}
\renewcommand{\baselinestretch}{1.00}
\AtBeginDocument{\fontsize{10.0pt}{12.0pt}\selectfont}
\setlength{\parskip}{4.0pt}
\setlength{\floatsep}{12.0pt plus 2pt minus 2pt}
\setlength{\textfloatsep}{16.0pt plus 2pt minus 2pt}
\setlength{\intextsep}{12.0pt plus 2pt minus 2pt}
\setlength{\abovecaptionskip}{6.0pt}
\setlength{\belowcaptionskip}{4.0pt}
\AtBeginDocument{\@ifundefined{bibsep}{}{\setlength{\bibsep}{3.0pt}}}
\widowpenalty=10000
\clubpenalty=10000
\makeatother
% J2L-OPT-END
\begin{document}
% >>> j2l compat: body-text size (+5%)
\makeatletter
\edef\JLTXbasesize{\f@size}
\@tempdimb=1.2\dimexpr\f@size pt\relax\edef\JLTXbaseskip{\strip@pt\@tempdimb}
\def\JLTXbodyscale{1.050}
\let\JLTX@normalsize\normalsize
\renewcommand\normalsize{%
  \JLTX@normalsize
  \@tempdima=\JLTXbodyscale\dimexpr\f@size pt\relax
  \@tempdimb=1.2\@tempdima
  \fontsize{\strip@pt\@tempdima}{\strip@pt\@tempdimb}\selectfont}%
\@ifundefined{@makecaption}{}{%
  \let\JLTX@makecaption\@makecaption
  \long\def\@makecaption#1#2{%
    \begingroup\fontsize{\JLTXbasesize}{\JLTXbaseskip}\selectfont
    \let\normalsize\JLTX@normalsize
    \JLTX@makecaption{#1}{#2}\endgroup}}%
\@ifundefined{thebibliography}{}{%
  \let\JLTX@thebibliography\thebibliography
  \def\thebibliography{%
    \fontsize{\JLTXbasesize}{\JLTXbaseskip}\selectfont\JLTX@thebibliography}}%
\makeatother
\normalsize
% <<< j2l compat


% Title block rendered inline (article's \maketitle forces a page break when
% preceded by other content, which pushed the whole paper to page 2).
\begin{center}
  {\LARGE\bfseries Performance Evaluation of MANET Protocols\par}
  \vspace{1em}
  {\normalsize Ali Researcher$^{1}$\par}
  \vspace{0.5em}
  {\small $^{1}$ Department of Networks, Test University\par}
\end{center}
\vspace{1em}

\begin{abstract}
This synthetic paper verifies extraction of native Office objects. Throughput trends are shown in Figure 1 and protocol share in Figure 2.
\end{abstract}

\noindent\textbf{Keywords:} charts, extraction, MANET
\vspace{1em}

\par\medskip\noindent {\large\bfseries Introduction\par}\smallskip

As Figure 1 demonstrates, throughput improves annually while latency drops. The distribution across protocols, given in Figure 2, favours AODV.

\begin{figure}[H]
\centering
\includegraphics[width=\linewidth,height=0.3\textheight,keepaspectratio]{media/office_obj_1.png}
\caption{Yearly throughput and latency of the testbed.}
\label{fig:office1}
\end{figure}

Beyond raw performance, the routing share matters for deployment planning as discussed above and revisited in Figure 2.

\begin{figure}[H]
\centering
\includegraphics[width=\linewidth,height=0.3\textheight,keepaspectratio]{media/office_obj_2.png}
\caption{Protocol share in surveyed deployments.}
\label{fig:office2}
\end{figure}

\par\medskip\noindent {\large\bfseries Discussion\par}\smallskip

Figure 3: Decision flow rendered as a native shape.

The shape in Figure 3 and the annotation box summarise the method. Refer back to Figure 1 for absolute numbers.

\begin{thebibliography}{99}
\bibitem{ref1}
A. Author, A Paper About Networks, Journal of Tests, 2024.

\end{thebibliography}


\end{document}
